\section*{Bab \ref{ch:b2:current-potential-curves} --- Kurva Arus--Potensial}

\begin{solution}{exo:b2:current-potential-curves:1}
Seng teroksidasi: \omterm{def:b2:current-potential-curves:ie-curve}{arus anodik}, positif. Ion-ion tembaga tereduksi pada
elektrode tembaga: \omterm{def:b2:current-potential-curves:ie-curve}{arus katodik}, negatif. Kedua arus itu sama besarnya,
karena muatan yang sama mengalir melalui rangkaian.
\end{solution}

\begin{solution}{exo:b2:current-potential-curves:2}
$|i_{\lim}| = 96\,485 \times 0.50 \times 10^{-4} \times 1.6 \times 10^{-9} \times
2.0/(25 \times 10^{-6}) = \qty{6.2e-4}{A}$, yaitu \qty{0.62}{mA}.
\end{solution}

\begin{solution}{exo:b2:current-potential-curves:3}
\omterm{def:b2:current-potential-curves:fast-slow}{Sistem cepat} memotong sumbu potensial secara tajam pada potensial
Nernst-nya; \omterm{def:b2:current-potential-curves:fast-slow}{sistem lambat} menunjukkan bentangan datar berarus nol di
sekitarnya, dengan cabang-cabang yang baru dimulai setelah sebuah
\omterm{def:b2:current-potential-curves:overpotential}{overpotensial}. Pada platina: \ce{Fe^3+}/\ce{Fe^2+} (cepat);
\ce{O2}/\ce{H2O} (lambat).
\end{solution}

\begin{solution}{exo:b2:current-potential-curves:4}
Pada potensial setengah gelombang, yang sama dengan
$E^\circ = \qty{0.77}{V}$ jika kedua bentuk berdifusi dengan cara yang
sama.
\end{solution}

\begin{solution}{exo:b2:current-potential-curves:5}
Di dekat \qty{0.77}{V}, gelombang katodik \ce{Fe^3+} (cepat), yang
mencapai plato sebanding dengan konsentrasinya. Di dekat \qty{0.34}{V},
gelombang kedua, ketika tembaga diendapkan; platonya menambah plato
pertama, sekitar dua kali lebih tinggi untuk konsentrasi yang sama dan
koefisien difusi yang serupa, karena $n = 2$. Lalu, di dekat \qty{0}{V}
pada pH 0 pada platina (lebih rendah setelah elektrode terlapisi tembaga),
dinding curam reduksi \ce{H+}.
\end{solution}

\begin{solution}{exo:b2:current-potential-curves:6}
Pada \qty{0.50}{V}: hanya \ce{Fe^3+} yang direduksi menjadi \ce{Fe^2+},
pada \omterm{def:b2:current-potential-curves:limiting-current}{arus pembatasnya}. Pada \qty{0.10}{V}: \ce{Fe^3+} direduksi dan tembaga
diendapkan bersamaan, masing-masing pada \omterm{def:b2:current-potential-curves:limiting-current}{arus pembatasnya}; totalnya adalah
jumlah kedua plato itu.
\end{solution}

\begin{solution}{exo:b2:current-potential-curves:7}
pH 0: \qty{0}{V} dan \qty{1.23}{V}; pH 14: \qty{-0.83}{V} dan
\qty{0.40}{V}. Pada platina, jendela itu mempertahankan lebarnya dan
bergeser ke bawah sebesar \qty{59}{mV} per satuan pH, dengan setiap
dinding mengikuti pasangannya, dan dinding anodik tetap tergeser ke atas
oleh \omterm{def:b2:current-potential-curves:overpotential}{overpotensial} oksigen.
\end{solution}

\begin{solution}{exo:b2:current-potential-curves:8}
Pada \qty{1.0}{V}, \ce{Fe^2+} dioksidasi pada platonya (\omterm{def:b2:current-potential-curves:ie-curve}{arus anodik}
sebanding dengan $[\ce{Fe^2+}]$) dan setiap \ce{Ce^4+} direduksi pada
platonya (\omterm{def:b2:current-potential-curves:ie-curve}{arus katodik} sebanding dengan $[\ce{Ce^4+}]$). Sebelum
ekuivalensi, \omterm{def:b2:current-potential-curves:ie-curve}{arus anodik} turun secara linear terhadap volume yang
ditambahkan; setelahnya, \omterm{def:b2:current-potential-curves:ie-curve}{arus katodik} bertambah secara linear. Kedua garis
lurus itu memotong nol pada ekuivalensi.
\end{solution}

\begin{solution}{exo:b2:current-potential-curves:9}
Plato-platonya menjadi dua kali lipat ($|i_{\lim}| \propto 1/\delta$).
Potensial setengah gelombang tidak berubah, karena lapisannya sama untuk
kedua bentuk; demikian pula potensial arus nol \omterm{def:b2:current-potential-curves:fast-slow}{sistem cepat}, yang
merupakan potensial Nernst.
\end{solution}

\begin{solution}{exo:b2:current-potential-curves:10}
Seperti dalam bab ini dengan $i_{\lim,c} = 0$: $E = E_{1/2} + (RT/nF)\ln[i/(i_{\lim} -
i)]$. Dalam logaritma desimal,
$E = E_{1/2} + (RT\ln 10/nF)\log[i/(i_{\lim} - i)]$, sebuah garis lurus
dengan kemiringan $0.0592/n$ volt pada \qty{298}{K}, yang memotong sumbu
$E$ pada $E_{1/2}$.
\end{solution}

\begin{solution}{exo:b2:current-potential-curves:11}
Jika sendirian, seng menetap pada potensial tempat \omterm{def:b2:current-potential-curves:ie-curve}{arus anodiknya}
(oksidasi seng yang cepat) sama dengan \omterm{def:b2:current-potential-curves:ie-curve}{arus katodik} reduksi \ce{H+} pada
seng, yang kecil karena reduksi itu lambat pada seng: pelarutannya lambat.
Jika menyentuh platina, elektron yang dilepaskan seng dapat mereduksi
\ce{H+} pada platina, tempat reduksi itu cepat: \omterm{def:b2:current-potential-curves:ie-curve}{arus katodik} pada potensial
tertentu jauh lebih besar, keseimbangannya tercapai pada arus yang lebih
tinggi, dan seng larut lebih cepat sementara hidrogen terbentuk pada
platina.
\end{solution}

\begin{solution}{exo:b2:current-potential-curves:12}
Pada platina, oksidasi air dimulai di dekat \qty{1.6}{V} pada pH 0:
sebelum \qty{2.0}{V}, arusnya akan habis untuk membuat dioksigen.
Elektrode tempat oksidasi air memerlukan \omterm{def:b2:current-potential-curves:overpotential}{overpotensial} yang lebih besar
mendorong dinding itu melampaui \qty{2.0}{V}, sehingga spesies itu dapat
dioksidasi di dalam jendela.
\end{solution}

\begin{solution}{pb:b2:current-potential-curves:1}
\textbf{1.} \ce{O2 + 2H2O + 4e- -> 4OH-}.
\textbf{2.} \ce{Ag + Cl- -> AgCl + e-} (empat kali); keseluruhan
\ce{O2 + 2H2O + 4Ag +
4Cl- -> 4AgCl + 4OH-}.
\textbf{3.} $E = 1.229 - 0.0592 \times 7 + (0.0592/4)\log 0.21 = \qty{0.80}{V}$.
\textbf{4.} $E^\circ(\ce{AgCl}/\ce{Ag}) = \qty{0.22}{V}$ dengan klorida
beraktivitas 1: $-0.60 + 0.22 = \qty{-0.38}{V}$ pada skala hidrogen.
\textbf{5.} Reduksi \ce{O2} adalah \omterm{def:b2:current-potential-curves:fast-slow}{sistem lambat}: di bawah potensial
Nernst-nya diperlukan \omterm{def:b2:current-potential-curves:overpotential}{overpotensial} beberapa persepuluh volt sebelum arus
bertambah, dan lebih banyak lagi untuk mencapai plato.
\textbf{6.} Pada plato, arus hanya bergantung pada pasokan \ce{O2}, bukan
pada pergeseran kecil potensial: arus itu mengukur konsentrasi.
\textbf{7.} Reduksi air menjadi dihidrogen (dinding katodik), yang akan
menambahkan arus yang tidak berhubungan dengan oksigen.
\textbf{8.} Linear melintasi membran, dari $c$ di sisi sampel sampai nol di
katode.
\textbf{9.} $j = D_m c/L$ per satuan luas.
\textbf{10.} $i = -4FAD_m c/L$ (katodik), $|i| = 4FAD_mc/L$.
\textbf{11.} Membran adalah tahap yang paling lambat: difusi melintasinya
jauh lebih lambat daripada di dalam air di luarnya, sehingga membran
berperan sebagai \omterm{def:b2:current-potential-curves:limiting-current}{lapisan difusi}.
\textbf{12.} Gradien konsentrasinya terletak di dalam membran, yang tebalnya
tidak bergantung pada pengadukan (asalkan air di luarnya terus
diperbarui); endapan pada membran menambah tebal dan menurunkan arus.
\textbf{13.} $0.2095 \times (101.325 - 3.17) = \qty{20.56}{kPa}$.
\textbf{14.} $c = 1.3 \times 10^{-5} \times 20\,560 = \qty{0.27}{mol/m^3}$ ($0.267$).
\textbf{15.} $0.267 \times 32.00 = \qty{8.6}{mg/L}$ (\unit{g/m^3}).
\textbf{16.} $|i| = 4 \times 96\,485 \times 1.0 \times 10^{-5} \times 1.0 \times 10^{-10}
\times 0.267/(25 \times 10^{-6}) = \qty{4.1}{\micro A}$.
\textbf{17.} $4.05/8.55 = \qty{0.474}{\micro A}$ per \unit{mg/L}.
\textbf{18.} Tebal membran dan koefisien difusinya tidak diketahui dengan
tepat serta berubah seiring umur dan suhu; pengukuran dalam larutan yang
diketahui mencakup semuanya.
\textbf{19.} Nol (dalam praktik ada arus sisa yang kecil, yang
dikurangkan).
\textbf{20.} $2.80/0.474 = \qty{5.9}{mg/L}$.
\textbf{21.} $2.80/4.05 = 69~\%$ kejenuhan.
\textbf{22.} Oksigen yang dikonsumsi oleh pembusukan bahan organik (limbah,
tumbuhan mati) lebih cepat daripada pasokannya dari udara dan dari
fotosintesis.
\textbf{23.} $2.80 \times 10^{-6}/(4 \times 96\,485) \times 3600 = \qty{2.6e-8}{mol}$
per jam: dapat diabaikan dibandingkan dengan oksigen dalam satu liter air
sekalipun (\qty{1.8e-4}{mol}).
\textbf{24.} Kelarutan \ce{O2} (\omterm{prop:b2:chemical-potential:henry}{tetapan Henry}) dan difusi melintasi membran
sama-sama berubah terhadap suhu: kepekaan yang diperoleh pada
\qty{25}{\degreeCelsius} tidak berlaku lagi.
\textbf{25.} Sampel itu mengandung $\boldsymbol{\approx \qty{5.9}{mg/L}}$
oksigen terlarut.
\end{solution}
